\section{Discussion}
In this report, we propose a variation to the coupling scheme introduced by Haas and Giles to reduce compute device memory challenges introduced by the proposed LUT solution \cite{haasNestedMLMCFramework2025a}. By doing so, we are able to explore and validate theoretical results described in their work using a GPU. While our limited capabilities for using varied precisions to analyse the effect of different bit-width optimisations on performance are not able to capture some of the specific goals of a nested multilevel scheme, we experimentally validate findings in CPU/GPU computing environments. 

These findings build directly upon the work done in MXB371, indicating potentially significant energy efficiency improvements for complex biochemical reaction networks by exploiting heterogeneous computing frameworks. As Systems Biology advances and the mathematical models required increase in complexity, we require these algorithms to ensure that wall-time on results of these models is reasonable. Importantly, simply adding more compute to the environment is not a sustainable solution to the problem, nor always feasible. Thus, the work in this report demonstrates methodologies that not only reduce the total computation time of the algorithms, but also the effective power usage. 

\subsection{Future Work}
There are several key avenues to pursue further work in this space. Before moving our implementations to accelerator hardware, i.e. FPGAs, we can continue to develop and validate methods on CPU/GPU architectures. Primarily, it would be of interest to compare the different methods proposed by Haas and Giles for coupling the precision levels. Further analysis on the effectiveness of our coupling scheme would also be insightful and how it compares to those discussed. In doing so, developing the capabilities to perform half precision, 16 bit, simulations on existing GPUs and coupling this to a double precision, 64 bit, path could allow us to better explore the effect of coupling. 

We further explore the benefits of the nested framework on traditional heterogeneous computing environments by scaling the already developed coupled kernels to be run in a High Performance Computing facility. This would allow us to scale our solution to further the power efficiency and wall time reductions explored in the \refsup{sec:supplementary_material}. Once this is developed, the shift to inference problems using the developed simulation algorithms could also be explored. 

The most obvious advancement of this work than would be to apply our algorithms on an FPGA and explore optimisations based on this hardware. By doing so, we are able to investigate more of the theoretical results proposed by Haas and Giles. Specifically, in optimising the bit-width, or word length as described in \cite{haasNestedMLMCFramework2025a}, for our context and understanding how this is impacted when applied to real hardware. Additionally, exploring how potentially multiple precision levels, similar to the multiple timestep levels, could increase performance further. 